% ch26.tex  --  Parallel Depth-First Search
% Transcribed from: Lec 26 -- Parallel Algorithms (Subodh Sharma, April 24, 2026)

% Give every ppbox a blank line between its label and its body.
\tcbset{ppbox/.append style={after title={\par\smallskip}}}

\chapter{Parallel Depth-First Search}

%------------------------------------------------------------------
% Learning outcomes
%------------------------------------------------------------------
\begin{ppkey}
After studying this chapter you should be able to:
\begin{itemize}
  \item Explain why DFS is intrinsically hard to parallelise and why a
        naive parallel expansion loses the depth-first property.
  \item Describe a depth biased work-pool parallel search and state what it does and
        does \emph{not} guarantee about the resulting tree.
  \item Outline the divide-and-conquer parallel DFS strategy: spine
        extraction, off-spine components, attachment, recursion, and
        merging.
  \item Apply parallel Kruskal's algorithm with Bor{\r u}vka rounds,
        including parallel deduplication, Union-Find with pointer
        jumping, and cross-component edge detection.
  \item Construct an Euler tour of a tree in parallel and use parallel
        list ranking via prefix sum to recover node depths.
  \item Identify off-spine components and attach each one to the correct spine
        vertex using a deterministic rule.
  \item Merge per-component DFS trees into one global DFS tree by
        computing global DFS numbers via parallel prefix sum.
  \item Classify non-tree edges as forward, back, or cross edges in
        $\mathcal{O}(1)$ per edge using DFS-number intervals.
    \item Understand the difference between parallel graph traversal and a verified DFS tree.
\end{itemize}
\end{ppkey}

%==================================================================
\section{Introduction}
%==================================================================
\begin{ppdefn}
A \ppterm{valid DFS tree} of a connected graph $G$ rooted at a vertex
$r$ is a rooted spanning tree $T$ of $G$ such that there exists a
sequential DFS execution starting from $r$ whose tree edges are exactly
the edges of $T$.
 
Equivalently, $T$ is a valid DFS tree if and only if every non-tree
edge of $G$ joins some node to one of its ancestors in $T$, that
is, every non-tree edge is a \emph{back edge}.  No non-tree edge of an
undirected graph may connect two nodes that are neither ancestor nor
descendant of each other.
% For directed graphs the definition admits two further non-tree edge
% types --- forward and cross edges --- and the classification carried
% out at the merge step (Section~\ref{sec:dfs-merge}) handles all three
% cases uniformly.
\end{ppdefn}
Depth-First Search is a workhorse of sequential graph algorithms: it
underpins biconnectivity, strongly connected components, topological
sort, and many shortest-path and matching algorithms.
Sequentially it is trivial; in parallel it is one of the harder
graph problems.
This chapter develops a parallel DFS that produces a valid DFS tree by
combining several primitives: parallel sort and
scan, parallel Union-Find with pointer jumping, Euler tours, and
parallel list ranking.  The treatment follows the textbook development
in~\cite{grama2003,kumar1994}.

Three obstacles make DFS resist parallelisation:
\begin{itemize}
  \item \textbf{Long sequential chains from depth exploration.}
        Each recursive call depends on the discovery state produced by
        the previous one, so the critical path is essentially a chain
        of node visits.
  \item \textbf{Constructing a valid DFS tree is hard.}
        A DFS tree is not just any spanning tree; it must be one that a
        sequential DFS could have produced.
  \item \textbf{Constructing a \emph{specific} DFS order is even harder.}
        The DFS stack can be shared across processors but orderings such as ``smallest numbered neighbour first'' will require the algorithm to know whether the smallest neighbour's subtree has finished. This creates a long dependence chain
\end{itemize}

\begin{ppexample}
Consider the small graph
\[
  A \to \{B, C\}, \qquad B \to \{D, E\}.
\]
A naive ``parallel DFS'' that expands all unvisited neighbours of $A$
simultaneously discovers $B$ and $C$ in the same step.
A sequential DFS, by contrast, would explore the entire subtree rooted
at $B$ (so $D$ and $E$) before ever touching $C$.
The parallel discovery pattern is no longer depth-first; it is just a
parallel traversal that happens to cover the graph.
We have given up the very property we wanted to preserve.
\end{ppexample}

\begin{figure}[ht]
  \centering
  \includegraphics[width=0.36\textwidth]{col380-book/chapters/figures/fig-naive-tree.jpg}
  \caption{The toy tree used in the naive-traversal argument.
           Expanding $B$ and $C$ in parallel from $A$ destroys the
           depth-first discovery order.}
  \label{fig:dfs-naive-tree}
\end{figure}

%==================================================================
\section{Work-Pool Parallel Search}
%==================================================================

Before attempting a ``proper'' parallel DFS, we record a practical
compromise: a depth-biased parallel search using a shared work pool.
The output is not guaranteed to be a strict DFS tree, but the
construction is simple, scalable, and a useful starting point for
applications that mainly need a spanning tree biased toward
depth~\cite{pacheco2011}.

\begin{ppdefn}
\ppterm{Work-pool parallel search} maintains a shared stack (or bag) of
frontier vertices.  Many processors repeatedly do the following:
\begin{enumerate}
  \item Remove a vertex $u$ from the bag.
  \item Inspect $u$'s adjacency list.
  \item For each unvisited neighbour $v$, atomically attempt to mark
        $v$ as visited.
  \item If the mark succeeds, set \texttt{parent[v] = u} and insert
        $v$ into the bag.
\end{enumerate}
\end{ppdefn}

Atomic marking ensures each vertex is claimed by exactly one processor,
so \texttt{parent[]} defines a valid spanning tree of the connected
component reached from the start vertex.

\begin{pppitfall}
A work-pool search returns a spanning tree, not a DFS tree.
The order in which vertices come out of the bag depends on processor
scheduling and contention, so a sequential DFS could not, in general,
have produced the same tree. The resultant parent tree may violate the DFS ancestor descendant property.
\end{pppitfall}

%==================================================================
\section{Divide-and-Conquer Parallel DFS}
\label{sec:dfs-divide-conquer}
%==================================================================

The remainder of the chapter develops a divide-and-conquer parallel DFS
that does produce a valid DFS tree.
The strategy is to identify a single root-to-leaf path of the DFS
tree, called the \ppterm{spine}, and then recurse independently on the
components that hang off it.

\begin{ppdefn}
The \ppterm{spine} of a DFS tree is a path from the root to a leaf
that we commit to in advance, before recursive subtrees are computed.
In the construction below, the spine is the longest root-to-leaf path
in a (minimum) spanning tree of the graph.
\end{ppdefn}

\begin{ppnote}
The high-level steps are:
\begin{enumerate}
  \item Identify a path that acts as the spine of the DFS tree.
  \item Find the connected components hanging off the spine.
  \item Assign each component to an attachment vertex on the spine.
  \item Recursively compute DFS trees of the components in parallel.
  \item Merge the recursive trees into one global DFS tree.
  \item Verify the DFS property, if needed.
\end{enumerate}
Steps 1 and 5 carry most of the algorithmic content; steps 2--4 are
bookkeeping built on parallel Union-Find.
\end{ppnote}

%==================================================================
\section{Step 1a: Computing the Spine via a Spanning Tree}
\label{sec:dfs-spanning}
%==================================================================

The spine is extracted from a (minimum) spanning tree of the input
graph.
We compute the spanning tree in parallel using a Bor{\r u}vka-style
algorithm; the lecture slides label it ``parallel Kruskal'' because
edges are processed in non-decreasing weight order, but the per-round
component contraction is pure Bor{\r u}vka~\cite{grama2003,kumar1994}.

\begin{figure}[ht]
  \centering
  \includegraphics[width=0.85\textwidth]{col380-book/chapters/figures/fig-input-graph.jpg}
  \caption{The running example used throughout the chapter.  Vertices
           $\{A, B, C, D, E, F\}$ with weighted undirected edges; the
           same graph is referenced by every subsequent figure.}
  \label{fig:dfs-input-graph}
\end{figure}

\subsection{Parallel Edge Sort}

Sort the edge list by weight in parallel using parallel merge sort.
Local chunks of size $E/p$ can be sorted in
$\mathcal{O}\!\bigl(\frac{E}{p}\log\frac{E}{p}\bigr)$ time, and the
merge reduction combines them in $\mathcal{O}(\log p)$ rounds.

\begin{ppexample}
Take the unsorted array $[38, 27, 43, 3, 9, 82, 10, 19]$ split across
two threads.  After local sort:
\begin{align*}
  T_1 : [38, 27, 43, 3]      &\;\longrightarrow\; [3, 27, 38, 43], \\
  T_2 : [9,\; 82,\; 10,\; 19] &\;\longrightarrow\; [9, 10, 19, 82].
\end{align*}
A parallel merge then combines the two sorted runs into one.
The merge step itself runs in parallel:
\begin{enumerate}
  \item Pick the median of the larger array $A$, say $A[\mathrm{mid}]$.
  \item Binary-search the position $k$ in $B$ where $A[\mathrm{mid}]$
        would fit.
  \item This yields four segments: the left/right halves of $A$ around
        $\mathrm{mid}$ and the left/right halves of $B$ around $k$.
  \item Concurrently, processor $P_1$ recursively merges the left
        halves and $P_2$ recursively merges the right halves.
\end{enumerate}
\end{ppexample}

\subsection{Bor{\r u}vka Rounds for Spanning-Tree Construction}

Initially every vertex is its own single-vertex component, owned by one
processor.  In each round we do the following in parallel:
\begin{itemize}
  \item Every component finds its minimum-weight outgoing edge.
  \item The selected edges are de-duplicated (two components can pick
        the same edge from opposite ends).
  \item Components are contracted into super-vertices.
\end{itemize}

\begin{ppexample}
Round~1 of Bor{\r u}vka on the graph from
Section~\ref{sec:dfs-spanning} produces:
\begin{center}
\begin{tabular}{@{}clll@{}}
\toprule
Proc.\ & Component & Min outgoing edge \\
\midrule
$P_1$ & $\{A\}$ & $(A,C),\; w=2$ \\
$P_2$ & $\{B\}$ & $(B,E),\; w=3$ \\
$P_3$ & $\{C\}$ & $(A,C),\; w=2$ \quad (duplicate of $P_1$) \\
$P_4$ & $\{D\}$ & $(B,D),\; w=6$ \\
$P_5$ & $\{E\}$ & $(B,E),\; w=3$ \quad (duplicate of $P_2$) \\
$P_6$ & $\{F\}$ & $(D,F),\; w=7$ \\
\bottomrule
\end{tabular}
\end{center}
Edges $(A,C)$ and $(B,E)$ are each selected by two processors and must
be removed before contraction.
\end{ppexample}

% (The table above already shows the per-component selections; slide 8
% in the lecture additionally highlights the duplicate selections in red.)

\subsection{Parallel Deduplication}

Duplicates are removed by sort + scan:
\begin{enumerate}
  \item Each undirected edge is normalised to a canonical form, e.g.\
        the endpoint with the smaller vertex id is written first.
  \item The list is sorted by this canonical key.
  \item An exclusive prefix scan flags the first occurrence of each
        edge; later copies are dropped.
\end{enumerate}
After deduplication on the running example we are left with
\[
  (A,C)\,w{=}2,\quad
  (B,E)\,w{=}3,\quad
  (B,D)\,w{=}6,\quad
  (D,F)\,w{=}7.
\]

\subsection{Parallel Union-Find with Pointer Jumping}

To contract the surviving edges into super-vertices we run a parallel
Union-Find.

\begin{ppdefn}
A \ppterm{Union-Find} (or disjoint-set) data structure exposes two
primitives:
\begin{itemize}
  \item \texttt{Find$(u)$} --- return the representative (root) of the
        set containing $u$.
  \item \texttt{Union$(u,v)$} --- merge the sets containing $u$ and $v$.
\end{itemize}
\end{ppdefn}

\begin{ppdefn}
\ppterm{Pointer jumping} is a parallel technique in which every node's
parent pointer is replaced by its \emph{grandparent}'s pointer:
\[
  \texttt{parent[v]} \;\leftarrow\; \texttt{parent[parent[v]]}.
\]
Applied in parallel by every processor and iterated until a fixed
point, pointer jumping flattens any rooted forest in $\mathcal{O}(\log n)$
rounds, after which one read of \texttt{parent[v]} returns the
root~\cite{kumar1994}.
\end{ppdefn}

\begin{ppnote}
The flattened Union-Find tree is used \emph{only} to identify roots; it
does not preserve the original graph or tree structure.
We will keep a separate \texttt{parent[]} array for the spanning tree
itself.
\end{ppnote}

After the first round of Union-Find on our four de-duplicated edges,
two components remain: $\{A,C\}$ rooted at $A$, and $\{B,D,E,F\}$
rooted at $B$.
The path from $F$ is compressed directly to $B$ during pointer jumping.

\begin{figure}[ht]
  \centering
  \includegraphics[width=0.85\textwidth]{col380-book/chapters/figures/fig-uf-round1.jpg}
  \caption{Two components after the first Union-Find round on the four
           de-duplicated edges.  Component $\{A,C\}$ is rooted at $A$;
           component $\{B,D,E,F\}$ is rooted at $B$.  The orange arrow
           from $F$ to $B$ is the path-compressed pointer produced by
           pointer jumping on edge $(D,F)$.}
  \label{fig:dfs-uf-round1}
\end{figure}

\subsection{Cross-Component Edge Detection}

After each round we still need to connect the surviving components.
For every edge $(u,v)$ in the input, one processor checks in $\mathcal{O}(1)$
whether the edge crosses components,
\[
  \texttt{Find}(u) \;\neq\; \texttt{Find}(v),
\]
and if so records $(u,v)$ together with its weight.
A parallel min-reduction over the collected list selects the next
minimum-weight inter-component edge.

\begin{ppexample}
For our running graph, the cross-edge selected in round 2 is $(A,B)$
with $w=4$, which connects $\{A,C\}$ to $\{B,D,E,F\}$.
After this round all six vertices belong to a single component and the
spanning tree is complete:
\[
  S = \{(A,C), (A,B), (B,E), (B,D), (D,F)\}.
\]
\end{ppexample}

\begin{figure}[ht]
  \centering
  \includegraphics[width=0.85\textwidth]{col380-book/chapters/figures/fig-spanning-tree.jpg}
  \caption{The spanning tree highlighted on the running graph.  Tree
           edges $\{(A,C),(A,B),(B,E),(B,D),(D,F)\}$ are drawn solid
           green; non-tree edges $\{(B,C),(C,D),(C,E),(E,F)\}$ are
           dashed.}
  \label{fig:dfs-spanning-tree}
\end{figure}

%==================================================================
\section{Step 1b: Extracting the Root-to-Leaf Spine}
\label{sec:dfs-spine}
%==================================================================

Given the spanning tree $S$, the spine is its longest root-to-leaf
path.  We extract it in two stages: an Euler tour to linearise the
tree, then parallel list ranking to recover depths.

\subsection{Euler Tour Construction}

\begin{ppdefn}
An \ppterm{Euler tour} of a tree replaces each undirected tree edge
$(u,v)$ with two directed arcs $u \to v$ and $v \to u$ and orders the
arcs so that they trace the order in which a sequential DFS would
follow them.  For each arc we compute a \texttt{next} pointer that
points to the arc following it in the tour~\cite{grama2003}.
\end{ppdefn}

The \texttt{next} pointers obey two rules:
\begin{itemize}
  \item \textbf{Rule 1 (descend).} $\texttt{next}(u \to v)$ is the
        first arc out of $v$ to $v$'s first child.
        If $v$ is a leaf then $\texttt{next}(u \to v) = (v \to u)$:
        the traversal returns immediately.
  \item \textbf{Rule 2 (backtrack).} $\texttt{next}(v \to u)$ is the
        next arc out of $u$ after $v$, that is, the arc to $u$'s next
        child following $v$.  If $v$ was the last child of $u$ then
        $\texttt{next}(v \to u) = (u \to \mathrm{parent}(u))$.
  \item The tour terminates when the return arc from the root's last
        child wraps back to the first arc out of the root.
\end{itemize}

This construction is purely local: each processor handles one arc and
writes one pointer, so the directed-graph build is an $\mathcal{O}(1)$
parallel computation per arc.

\begin{ppexample}
For the spanning tree of our running example, ordering children by
ascending label and starting from $A$, the Euler tour produces the arc
sequence
\[
\begin{aligned}
A{\to}B,\; B{\to}D,\; D{\to}F,\; F{\to}D,\; D{\to}B,\\
B{\to}E,\; E{\to}B,\; B{\to}A,\; A{\to}C,\; C{\to}A.
\end{aligned}
\]
Down-arcs descend into a new child; up-arcs return to a parent.
\end{ppexample}

\begin{figure}[ht]
  \centering
  \includegraphics[width=0.95\textwidth]{col380-book/chapters/figures/fig-euler-tour.jpg}
  \caption{Euler tour of the spanning tree as an ordered list of
           directed arcs (positions $0$--$9$).  Down-arcs (green) descend
           into a new child; up-arcs (purple) return to a parent.}
  \label{fig:dfs-euler-tour}
\end{figure}

\subsection{Parallel List Ranking via Prefix Sum}

Each arc carries a contribution: $+1$ for a down-arc and $0$ for an
up-arc.
A parallel exclusive prefix sum over these contributions assigns every
arc its position along a depth-counted ranking.
Reading the down-arcs gives node depths, and the maximum depth
identifies the deepest leaf.

\begin{ppexample}
On our running example, the prefix sum reveals that $F$ is the deepest
leaf at depth~3.  Following parent pointers from $F$ back to the root
recovers the spine
\[
  P \;=\; A \to B \to D \to F.
\]
Parent pointers can themselves be jumped in $\mathcal{O}(\log n)$ rounds, so
the recovery is also parallel.
\end{ppexample}

% (Slide 15 in the lecture provides a two-row tabular view of the
% list-ranking computation: the top row shows per-arc accumulators
% and \texttt{next} pointers, the bottom row shows the prefix-sum
% totals 4, 10, 17, 17, 17, 20, 20, 20, 22, 22.  The same data is
% reproduced in Exercise 5.)

\begin{ppnote}
When components were merged during Bor{\r u}vka contraction, depths of
nodes inside an absorbed component were incremented to reflect their
new position in the larger tree.
This depth bookkeeping is what allows the maximum-depth scan to identify
$F$ as the deepest leaf in a single global pass.
\end{ppnote}

%==================================================================
\section{Step 2: Components Hanging Off the Spine}
%==================================================================

With the spine $P = A \to B \to D \to F$ in hand, contract (or
mark-as-removed) the spine vertices and find the connected components
of the remaining off-spine sub-induced graph.
The same parallel Union-Find with pointer jumping is used, but on an
auxiliary parent relation.

For each off-spine node $v$ we initialise
\[
  \texttt{off\_parent[v]} \;=\;
  \begin{cases}
    \texttt{parent[v]} & \text{if } \texttt{parent[v]} \text{ is also off-spine},\\
    v                  & \text{otherwise (i.e.\ parent is on the spine)}.
  \end{cases}
\]
Then in each round every processor simultaneously executes
\[
  \texttt{off\_parent[v]} \;\leftarrow\;
    \texttt{off\_parent[\,off\_parent[v]\,]}.
\]
After $\mathcal{O}(\log n)$ rounds the auxiliary forest is flat, and every
off-spine node knows its component root.

\begin{ppexample}
On the running graph, two singleton components emerge:
$\mathrm{comp}_1 = \{C\}$ and $\mathrm{comp}_2 = \{E\}$.
Their attaching cross-edges to the spine are $(A,C)$ and $(B,E)$
respectively.
\end{ppexample}

\begin{figure}[ht]
  \centering
  \includegraphics[width=0.85\textwidth]{col380-book/chapters/figures/fig-spine-components.jpg}
  \caption{Spine $A \to B \to D \to F$ (orange) and the two off-spine
           components $\mathrm{comp}_1 = \{C\}$ hanging off $A$ and
           $\mathrm{comp}_2 = \{E\}$ hanging off $B$ (green dashed
           edges).  Spine depths $d = 0, 1, 2, 3$ are labelled above
           each spine vertex.}
  \label{fig:dfs-spine-components}
\end{figure}

%==================================================================
\section{Step 3: Component-to-Spine Attachment}
%==================================================================

A component may have edges to several spine vertices, so a deterministic
rule must pick a unique attachment.

\begin{ppdefn}
We attach component $C_i$ to the \ppterm{deepest} spine vertex it
connects to.  Algorithmically, we inspect all edges $(u, v_j)$ with
$u \in C_i$ and $v_j \in P$, compute the maximum index $j$ in
parallel, and attach $C_i$ to $v_j$.
\end{ppdefn}

\begin{ppnote}
Attaching to the deepest reachable spine vertex (rather than the
shallowest) keeps the merged DFS tree consistent with a sequential DFS
that would have explored as deeply as possible before backtracking.
A shallower attachment would force back-edges to become tree edges and
violate DFS-tree invariants.
\end{ppnote}

%==================================================================
\section{Step 4: Recursive DFS on Components}
%==================================================================

Each component $C_i$ rooted at its first entry vertex $\mathrm{root}_i$
becomes an independent recursive sub-problem.  We invoke
\texttt{DFS$(C_i, \mathrm{root}_i)$} on each one, and the recursive
calls run in parallel because the components are disjoint by
construction.

\begin{ppnote}
$\mathrm{root}_i$ can be picked deterministically: either the endpoint
of the first edge in $C_i$'s adjacency list, or the smallest-labelled
endpoint of the cross-edge that attaches $C_i$ to the spine.
Either choice fixes a unique recursive root and is reproducible across
runs.
\end{ppnote}

%==================================================================
\section{Step 5: Parallel Tree Merging}
\label{sec:dfs-merge}
%==================================================================

The recursion produces, for each $C_i$, a local DFS tree $T_i$ with its
own local DFS numbering.  We now stitch them onto the spine to form one
global DFS tree.

\subsection{Merge Invariants}

\begin{ppdefn}
A correct DFS-tree merge preserves three invariants:
\begin{description}
  \item[I1.] Each attachment edge
        $(\text{spine\_node} \to C_i\text{\_root})$ becomes a tree
        edge of the merged DFS tree.
  \item[I2.] DFS numbering respects pre-order: every spine node $s$
        is numbered before all nodes in the components hanging off
        $s$.
  \item[I3.] Non-tree edges are classified as forward, back, or cross
        edges using the DFS-number intervals of the merged tree.
\end{description}
\end{ppdefn}

\subsection{Computing Global DFS Numbers}

For each spine vertex $s_k$ we want:
\begin{itemize}
  \item the size of the component $C_i$ hanging off $s_k$, computed by
        a parallel reduction in $\mathcal{O}(\log^2 n)$;
  \item the global DFS number of $s_k$, given by
        \[
          \mathrm{dfs}[s_k]
            \;=\; k \;+\; \sum_{j < k} \mathrm{block\_size}[s_j].
        \]
\end{itemize}

The trick is to build an interleaved sequence whose slots alternate
``spine node / hanging block''
\[
  \underbrace{s_0}_{1}\;
  \underbrace{\text{block}(s_0)}_{|C_0|}\;
  \underbrace{s_1}_{1}\;
  \underbrace{\text{block}(s_1)}_{|C_1|}\;
  \cdots
\]
Each spine slot contributes $1$, and each block slot contributes its
component size (or $0$ if no component hangs off that spine vertex).
A parallel exclusive prefix sum over this contribution array yields each
slot's global DFS start number.

\begin{ppexample}
On our running graph the contribution array is
$[1,1,1,1,1,0,1,0]$, with the zeros corresponding to the empty blocks
under spine vertices $D$ and $F$.
The prefix sum produces the merged pre-order
\[
  A_0,\quad C_1,\quad B_2,\quad E_3,\quad D_4,\quad F_5,
\]
where subscripts denote global DFS numbers.
This satisfies invariants I1 and I2 by construction.
\end{ppexample}

\begin{figure}[ht]
  \centering
  \includegraphics[width=0.7\textwidth]{col380-book/chapters/figures/fig-merge-preorder.jpg}
  \caption{Merged DFS pre-order produced by the interleaved prefix sum
           on the contribution array $[1,1,1,1,1,0,1,0]$.  Spine
           vertices (orange) and off-spine vertices (green) are laid out
           in the order they receive their global DFS numbers.}
  \label{fig:dfs-merge-preorder}
\end{figure}

\subsection{Classifying Non-Tree Edges}

Once every node has a global DFS number $\mathrm{dfs}[\cdot]$ and a
subtree size $\mathrm{size}[\cdot]$, every non-tree edge $(u, v)$ in
the original graph is classified in parallel and in $\mathcal{O}(1)$ per
edge:

\begin{ppdefn}
For a non-tree edge $(u, v)$ in the merged DFS tree:
\begin{description}
  \item[Forward edge.]
        $\mathrm{dfs}[u] < \mathrm{dfs}[v]$ \emph{and}
        $\mathrm{dfs}[v] \le \mathrm{dfs}[u] + \mathrm{size}[u] - 1$.
        ($v$ lies in $u$'s subtree.)
  \item[Back edge.]
        $\mathrm{dfs}[v] < \mathrm{dfs}[u]$ \emph{and}
        $\mathrm{dfs}[u] \le \mathrm{dfs}[v] + \mathrm{size}[v] - 1$.
        ($u$ is a descendant of $v$.)
  \item[Cross edge.] Otherwise.
\end{description}
\end{ppdefn}

This is invariant I3, expressed as a constant-time predicate over
DFS-number intervals.
Because every edge is classified independently, the whole non-tree edge
classification is a single parallel pass.

%==================================================================
\section*{Key Takeaways}
\addcontentsline{toc}{section}{Key Takeaways}
%==================================================================

\begin{ppnote}
\begin{itemize}
  \item \textbf{Naive parallel DFS fails} because expanding all unvisited
        neighbours simultaneously discards the depth-first ordering we
        wanted in the first place.
  \item \textbf{Work-pool search} gives a depth-biased spanning tree
        cheaply but does not produce a valid DFS tree.
  \item \textbf{Divide-and-conquer DFS} extracts a spine, recurses on
        off-spine components in parallel, and merges results.
  \item \textbf{Parallel Kruskal/Bor{\r u}vka} builds the spanning tree
        by repeatedly picking minimum-weight outgoing edges per
        component, deduplicating with a parallel scan, and contracting
        with parallel Union-Find.
  \item \textbf{Pointer jumping} flattens any rooted forest in
        $\mathcal{O}(\log n)$ rounds: $\texttt{parent[v]} \leftarrow
        \texttt{parent[parent[v]]}$.
  \item \textbf{Euler tour + list ranking} converts a tree into a
        linear sequence with depth labels using only an
        $\mathcal{O}(1)$-per-arc construction and one parallel prefix sum.
  \item \textbf{The deepest-vertex attachment rule} is what keeps the
        merged tree consistent with a sequential DFS; shallower
        attachments would violate DFS-tree invariants.
  \item \textbf{Tree merging} reduces to a single interleaved prefix sum
        over $[\,1, |C_0|, 1, |C_1|, \ldots\,]$ that yields global DFS
        numbers.
  \item \textbf{Edge classification} is then $\mathcal{O}(1)$ per non-tree
        edge using DFS-number intervals: forward, back, or cross.
\end{itemize}
\end{ppnote}

%==================================================================
\section{Exercises}
%==================================================================

\begin{ppexercises}

\ppexercise{%
\textbf{Why naive parallelism fails.}  Take the small tree
$A \to \{B, C\}$ and $B \to \{D, E\}$ with the convention that
sequential DFS visits the smaller-labelled child first.
\begin{enumerate}[label=(\alph*)]
  \item List the discovery order produced by sequential DFS.
  \item List a possible discovery order if all unvisited neighbours of
        the current frontier are expanded simultaneously.
  \item Identify one tree edge produced by the naive parallel scheme
        that could not appear in any valid DFS tree.
\end{enumerate}

\textbf{Solution.}
\begin{enumerate}[label=(\alph*)]
  \item Sequential DFS: $A, B, D, E, C$.
  \item Naive parallel: $A$, then $\{B, C\}$ together, then
        $\{D, E\}$ together --- discovery order
        $A, B, C, D, E$ (or with $B, C$ swapped).
  \item In sequential DFS at $A$, after entering $B$ we must finish the
        $B$-subtree before exploring $C$; therefore $C$ is discovered
        from $A$ \emph{after} $D$ and $E$.  The naive parallel scheme
        discovers $C$ at depth~1 \emph{before} $D$ and $E$, an
        ordering no sequential DFS could produce.
\end{enumerate}
}

\ppexercise{%
\textbf{Tracing Bor{\r u}vka rounds.}  On the graph of
Section~\ref{sec:dfs-spanning}, perform the algorithm round by round.
\begin{enumerate}[label=(\alph*)]
  \item List each component's minimum-weight outgoing edge in round 1.
  \item Identify the duplicate selections.
  \item Show the contracted super-vertices after round 1.
  \item Find the next cross-component edge selected in round 2.
\end{enumerate}

\textbf{Solution.}
\begin{enumerate}[label=(\alph*)]
  \item $\{A\}{:}\,(A,C)\,2$,\;
        $\{B\}{:}\,(B,E)\,3$,\;
        $\{C\}{:}\,(A,C)\,2$,\;
        $\{D\}{:}\,(B,D)\,6$,\;
        $\{E\}{:}\,(B,E)\,3$,\;
        $\{F\}{:}\,(D,F)\,7$.
  \item $(A,C)$ and $(B,E)$ each appear twice.
  \item Two super-vertices: $\{A,C\}$ and $\{B,D,E,F\}$.
  \item Among edges crossing these two components --- $(A,B)\,4$,
        $(B,C)\,5$, $(C,D)\,8$, $(C,E)\,10$ --- the minimum is
        $(A,B)$ with $w=4$.
\end{enumerate}
}

\ppexercise{%
\textbf{Pointer jumping.}  Suppose Union-Find currently has
\[
  \texttt{parent} = [A,A,B,A,B,B], \qquad
  \text{indices } [A,B,C,D,E,F],
\]
where $\texttt{parent[X]}$ is the parent of $X$ (so $\texttt{parent[A]}=A$
means $A$ is a root).  Apply pointer jumping until a fixed point.
\begin{enumerate}[label=(\alph*)]
  \item State the parent array after one parallel round.
  \item State the parent array after the round at which a fixed point
        is reached.
  \item Why does the fixed point identify the root of each component?
\end{enumerate}

\textbf{Solution.}
The initial array is already shallow: every non-root has its grandparent
equal to its parent (every parent is itself a root).
\begin{enumerate}[label=(\alph*)]
  \item After one round, $\texttt{parent}$ is unchanged --- a fixed
        point is reached immediately.
  \item Same as (a): $[A,A,B,A,B,B]$.
  \item At the fixed point, $\texttt{parent[parent[v]]} = \texttt{parent[v]}$
        for every $v$, which means $\texttt{parent[v]}$ is its own root.
        A single read of $\texttt{parent[v]}$ now returns the component
        root in $\mathcal{O}(1)$.
\end{enumerate}
}

\ppexercise{%
\textbf{Euler tour next pointers.}  Take the rooted tree on $\{A,B,C\}$
where $A$ is the root with children $B$ and $C$ (in that order).
\begin{enumerate}[label=(\alph*)]
  \item List the four directed arcs of the Euler tour.
  \item For each arc, give its $\texttt{next}$ pointer using Rules 1
        and 2.
  \item Verify that the tour terminates correctly when the return arc
        from the last child wraps back to the first arc out of the
        root.
\end{enumerate}

\textbf{Solution.}
\begin{enumerate}[label=(\alph*)]
  \item Arcs: $A{\to}B$, $B{\to}A$, $A{\to}C$, $C{\to}A$.
  \item Pointers:
        \begin{itemize}
          \item $\texttt{next}(A{\to}B)$: $B$ is a leaf, so by Rule~1
                it is $B{\to}A$.
          \item $\texttt{next}(B{\to}A)$: $B$ is not the last child of
                $A$ (next child is $C$), so by Rule~2 it is $A{\to}C$.
          \item $\texttt{next}(A{\to}C)$: $C$ is a leaf, so by Rule~1
                it is $C{\to}A$.
          \item $\texttt{next}(C{\to}A)$: $C$ is the last child of the
                root, so by Rule~2 the tour terminates.
        \end{itemize}
  \item The arc sequence is $A{\to}B,\,B{\to}A,\,A{\to}C,\,C{\to}A$.
        After $C{\to}A$ there is no further arc, which is the
        termination condition for the root's last child.
\end{enumerate}
}

\ppexercise{%
\textbf{List ranking for depth.}  Suppose an Euler tour is encoded as a
sequence of contributions $[1, 1, 1, 0, 0, 1, 0, 0, 1, 0]$ (1 for a
down-arc, 0 for an up-arc).  Assume each down-arc enters a node whose
label appears in order $B, D, F, \cdot, \cdot, E, \cdot, \cdot, C,
\cdot$ (dots are up-arcs).
\begin{enumerate}[label=(\alph*)]
  \item Apply an exclusive prefix sum to obtain depth labels.
  \item Identify the deepest leaf and its depth.
  \item State the recovered root-to-leaf path.
\end{enumerate}

\textbf{Solution.}
\begin{enumerate}[label=(\alph*)]
  \item Exclusive prefix sum:
        $[0, 1, 2, 3, 3, 3, 4, 4, 4, 5]$.
        The values at down-arc positions give the discovery depth
        (offset by 1) of each new node.
  \item The maximum depth among down-arcs is $3$, attained at the
        down-arc that discovers $F$.  So $F$ is the deepest leaf at
        depth $3$.
  \item Recovering parents from $F$ gives $F \to D \to B \to A$, so
        the spine is $A \to B \to D \to F$.
\end{enumerate}
}

\ppexercise{%
\textbf{Off-spine components via \texttt{off\_parent}.}  Suppose the
spanning tree has $\texttt{parent}[X]$ as follows:
$\texttt{parent}[B]=A$, $\texttt{parent}[C]=A$, $\texttt{parent}[D]=B$,
$\texttt{parent}[E]=B$, $\texttt{parent}[F]=D$, with $A$ as the root.
The spine is $\{A, B, D, F\}$.
\begin{enumerate}[label=(\alph*)]
  \item Initialise $\texttt{off\_parent}$ for every off-spine node.
  \item Apply one round of pointer jumping.
  \item State the off-spine components.
\end{enumerate}

\textbf{Solution.}
The off-spine nodes are $\{C, E\}$.

\begin{enumerate}[label=(\alph*)]

\item \texttt{parent[C]} = $A$ is on the spine, so
\[
\texttt{off\_parent[C]} = C.
\]
\texttt{parent[E]} = $B$ is on the spine, so
\[
\texttt{off\_parent[E]} = E.
\]

\item One round:
\[
\begin{aligned}
\texttt{off\_parent[C]} 
&= \texttt{off\_parent[off\_parent[C]]} \\
&= \texttt{off\_parent[C]} \\
&= C.
\end{aligned}
\]
Similarly for $E$. Already at a fixed point.

\item Each off-spine node is its own root, giving two singleton
components $\{C\}$ and $\{E\}$.

\end{enumerate}
}

\ppexercise{%
\textbf{Edge classification by intervals.}  After merging, suppose
\[
\begin{aligned}
\mathrm{dfs} &= \{A{:}0, C{:}1, B{:}2, E{:}3, D{:}4, F{:}5\}, \\
\mathrm{size} &= \{A{:}6, C{:}1, B{:}4, E{:}1, D{:}2, F{:}1\}.
\end{aligned}
\]
Classify the following non-tree edges as forward, back, or cross:
$(A, D)$, $(D, A)$, $(C, E)$, $(E, F)$.

\textbf{Solution.}
\begin{itemize}
  \item $(A,D)$: $\mathrm{dfs}[A]=0<\mathrm{dfs}[D]=4$ and
        $4 \le 0 + 6 - 1 = 5$.  \textbf{Forward.}
  \item $(D,A)$: $\mathrm{dfs}[A]=0<\mathrm{dfs}[D]=4$ and
        $4 \le 0 + 6 - 1 = 5$.  \textbf{Back} ($D$ is a descendant of $A$).
  \item $(C,E)$: $\mathrm{dfs}[C]=1<\mathrm{dfs}[E]=3$ but
        $3 \not\le 1 + 1 - 1 = 1$, so not forward.  Symmetric check
        also fails.  \textbf{Cross.}
  \item $(E,F)$: $\mathrm{dfs}[E]=3<\mathrm{dfs}[F]=5$ but
        $5 \not\le 3 + 1 - 1 = 3$, so not forward.  Symmetric check
        also fails.  \textbf{Cross.}
\end{itemize}
}

\ppexercise{%
\textbf{Why deepest, not shallowest?}  Suppose component $C_i$ has
edges to spine vertices $s_1$ (depth 1) and $s_3$ (depth 3).
\begin{enumerate}[label=(\alph*)]
  \item Argue informally that attaching $C_i$ to $s_3$ is consistent
        with a sequential DFS that has already descended to $s_3$
        before recursing into $C_i$.
  \item Show that attaching $C_i$ to $s_1$ would force the edge
        $(s_3, C_i)$ to violate the DFS-tree invariant in
        Section~\ref{sec:dfs-merge}.
\end{enumerate}

\textbf{Solution.}
\begin{enumerate}[label=(\alph*)]
  \item A sequential DFS from the root walks down $s_0, s_1, s_2, s_3$
        before backtracking.  At $s_3$ it sees an edge into $C_i$ and
        descends into the component.  When it later climbs back through
        $s_2, s_1, s_0$, the edge $(s_1, C_i)$ is already a back-edge
        from a node already visited, not a new tree edge.  Attaching
        $C_i$ at $s_3$ matches this picture.
  \item If we instead attached $C_i$ at $s_1$, then $(s_1, C_i)$ would
        be the tree edge.  But during the sequential DFS, $s_3$ is
        already on the spine and $C_i$ is reachable from $s_3$ before
        we ever return to $s_1$.  The edge $(s_3, C_i)$ would have to
        be a forward edge from $s_3$ \emph{into} $C_i$'s subtree, which
        contradicts I3 because $C_i$'s DFS interval would lie outside
        $s_3$'s subtree interval.
\end{enumerate}
}

\end{ppexercises}

%==================================================================
\section*{References}
\addcontentsline{toc}{section}{References}
%==================================================================

The following sources were used in preparing this chapter.  Each is
cited at the points in the chapter where its treatment was drawn on
most directly.

\begin{itemize}
  \item A.~Grama, G.~Karypis, V.~Kumar, and A.~Gupta,
        \emph{Introduction to Parallel Computing}, 2nd~ed.
        Addison-Wesley, 2003.~\cite{grama2003}
        \\ \emph{Used for:} the parallel
        Bor{\r u}vka/Kruskal spanning-tree construction
        (Section~\ref{sec:dfs-spanning}) and the Euler-tour technique
        for tree linearisation (Section~\ref{sec:dfs-spine}).
  \item V.~Kumar, A.~Grama, A.~Gupta, and G.~Karypis,
        \emph{Introduction to Parallel Computing: Design and Analysis
        of Algorithms}.
        Benjamin/Cummings, 1994.~\cite{kumar1994}
        \\ \emph{Used for:} the divide-and-conquer parallel-DFS
        framework, parallel Union-Find with pointer jumping, and
        parallel list ranking via prefix sum.
  \item P.~Pacheco, \emph{An Introduction to Parallel Programming}.
        Morgan Kaufmann, 2011.~\cite{pacheco2011}
        \\ \emph{Used for:} the practical work-pool/shared-bag
        formulation of depth-biased parallel search in
        Section~2 (``Work-Pool Parallel Search'').
\end{itemize}

%==================================================================
\section*{Contributor}
\addcontentsline{toc}{section}{Contributor}
%==================================================================

\begin{ppnote}
Transcribed by \textbf{Priyanka Sar}, Entry No.\ \textbf{2023CS50578}
\end{ppnote}